\section{Hidden Rationale Queries (L4)}

\subsection{Overview}
\MFUsentencecase{\forthlevel} are the most challenging type of queries to address. Unlike \thirdlevel , which provide clear guidance on the rationales needed to respond to queries, \forthlevel involve domain-specific reasoning method that may not be explicitly described and are too numerous to exhaust. These rationales often encompass a wide variety that cannot be fully explored within the typical context window and may lack clear instructions, representing a form of domain expertise that is implicit within the data. Such data might include, but is not limited to:

\begin{itemize}
    \item \textbf{In-domain Data}: \MFUsentencecase{\forthlevel} might utilize data from the same domain, such as historical question-and-answer records or artificially generated data. This in-domain data inherently contains the reasoning skills or methodologies necessary to address current queries. For instance, in the context of Python programming puzzles, solutions to historical problems often include classical algorithms and problem-solving strategies that could aid in resolving current issues.
    \item \textbf{Preliminary Knowledge}: Another form of hidden rationales consists of extensive, dispersed knowledge bases that vary in application across different scenarios. This preliminary knowledge may constitute a comprehensive axiomatic system, such as all local legal codes that form the basis for legal judgments. It could also include proven intermediate conclusions that simplify reasoning processes in fields like mathematical proofs. When addressing real-world issues using external data, this prior knowledge might also stem from complex accumulations of human experiences and empirical summaries.
\end{itemize}

Navigating \forthlevel thus demands sophisticated analytical techniques to decode and leverage the latent wisdom embedded within disparate data sources, presenting significant challenges for RAG systems in effectively interpreting and applying such intricate and implicit information.

Here are some examples of queries at this level:
\begin{itemize}
    \item \textit{How will the economic situation affect the company's future development?} (given a collection of financial reports, with economic and financial rationale required)
    \item \textit{How to achieve 24 points using the numbers 5, 5, 5, and 1?} (given a series of 24-point game examples and corresponding answers.)
    \item \textit{Does Afghanistan permit a parent to confer his or her citizenship on a child born abroad?} (given the GLOBALCIT citizenship law dataset~\cite{guha2024legalbench})
\end{itemize}

\subsection{Challenges and Solutions}

The construction of data-augmented LLM applications is significantly challenged by \forthlevel, with primary difficulties manifesting in the following areas:

\begin{itemize}
    \item \textbf{Logical retrieval}: For questions involving hidden rationales, the helpfulness of external data does not simply depend on entity-level or semantic similarity, but rather on logical congruence or thematic alignment. Standard retrieval methods often struggle to capture the true target of the query or to identify text segments with logical similarities based on the problem presented. This necessitates the development of more sophisticated retrieval algorithms that can parse and identify underlying logical structures rather than relying solely on superficial textual similarities.
    \item  \textbf{Data insufficiency}: Fundamentally, external data may not explicitly contain the guidance or answers relevant to the current query. Instead, relevant information is often embedded in dispersed knowledges or illustrated through examples. This indirect presentation demands robust capabilities in data interpretation and synthesis, requiring LLMs to effectively derive coherent answers from fragmented or tangentially related data sources. Such challenges underscore the imperative for sophisticated data integration and reasoning capabilities within LLM frameworks to navigate the complexities of \forthlevel effectively.
\end{itemize}

\subsection{Offline Learning}
To address these types of queries, a common approach is to identify and extract rules and guidelines from datasets offline, subsequently retrieving related items. For generate reasoning rationales, some work like STaR~\cite{zelikman2022star} and LXS~\cite{stammer2023learning} used LLM for rationale generation. The former employs an iterative few-shot example method to generate from small dataset to large dataset, the later introduced a two-role explaining extraction process where a learner model generated explanations and a critic model assess them for validation. 
% \comment{LXS mainly for fine-tuning?}

GL~\cite{pang2023guideline} identifies errors and generalizes them into guidelines for future tasks through an in-context-learning. LEAP~\cite{zhang2024context} forms principles by generating mistakes, low-level principles, and high-level principles, incorporating these principles into prompts for final inference. RICP~\cite{sun2024retrieved} uses mistakes from training data to generate high-level reasoning and specific insights, then employs hierarchical clustering to group modes of errors, generating task-level and question-level principles, which are combined and retrieved for question-level insights. A Buffer-of-Thought~\cite{yang2024buffer} uses a problem distiller to distill a meta-buffer across many reasoning tasks.

Some integrated methods, like MedPrompt~\cite{nori2023can}, include GPT-4-generated chains of thought for training examples with self-validation, using these in conjunction with a KNN retrieval in-context learning approach. Agent Hospital~\cite{li2024agent} generates rationales through reflection and utilizes both record retrieval and experience retrieval on generated data.

Although these concepts go by many different names—such as guidelines, principles, experiences, and thought templates—the main idea is to extract common useful rationales to enhance reasoning queries. These rationales may come from self-generated chains of thought (MedPrompt, Buffer-of-Thought), training set mistakes (GL, RICP, Agent Hospital), or intentionally generated mistakes (LEAP). Additionally, some principles are used across all tasks (Agent Hospital, RICP), while others are dynamically retrieved for specific questions (MedPrompt, Buffer-of-Thought). Many of these works demonstrate that learning from cases to accumulate experience as rationales is beneficial for various reasoning tasks.

\subsection{In Context Learning (ICL)}

Using examples for in-context learning is a common method for uncovering hidden rationales. Pre-trained large language models exhibit substantial in-context learning capabilities, which can be enhanced by retrieving examples based on similarity, thereby leveraging the models' few-shot learning abilities~\cite{liu2021makes, brown2020language}. However, the inclusion of irrelevant information in prompts can easily distract LLMs, leading to incorrect responses~\cite{shi2023large,chen2023many}. OpenICL, as developed by Wu et al.~\cite{wu2023openicl}, constructed an ICL framework that explores the impact of different traditional methods of retrieving examples and inference techniques on ICL effectiveness.

Furthermore, training smaller models based on the feedback from LLMs regarding context examples to select optimal demonstrations and examples can improve the construction of context for specific tasks in a more targeted manner~\cite{li2023unified, wang2024large, ye2023compositional}. To address the issue that semantic similarity-based example retrieval may not cover the broader range of associations needed in practical tests, Su et al.\cite{su2022selective} employed an unsupervised, graph-based selective annotation method called vote-k, constructing a more diverse and representative example database for few-shot learning. Additionally, Zhang et al.\cite{zhang2022automatic} proposed an Auto-CoT method that clusters examples into various representative types. By sampling problems diversely and generating reasoning chains, this method constructs examples that better support the learning process.

However, enabling LLMs to master reasoning capabilities outside their trained domains through few-shot learning remains a substantial challenge. Wang et al. addressed this by sampling a variety of reasoning paths and marginalizing over these paths to select the most consistent answer, thereby enhancing the probability of LLMs selecting correct reasoning chains~\cite{wang2022self}. Agarwal et al. introduced two scalable methods for generating usable example, namely reinforced ICL and unsupervised ICL, that aim to replace human-generated examples, thus expanding the pool of available examples~\cite{agarwal2024many}. DIN-SQL~\cite{pourreza2024din} sought to decompose tasks into simpler subtasks and used the solutions to these sub-problems as prompts for LLMs, significantly improving their performance in generating SQL from text. Similarly, DUP~\cite{zhong2024achieving} identified three main issues LLMs face when using the chain of thought approach to solve complex mathematical word problems: semantic misunderstandings, computational errors, and missing steps, with semantic misunderstandings being a primary limiting factor. Encouraging LLMs to deeply understand the problems and extract essential information for resolution can significantly enhance their ability to solve mathematical problems by addressing these semantic misunderstandings.

In-context learning is increasingly being utilized across various fields such as mathematics, law, medicine, and finance~\cite{schumacher2024team, dou2023plugmed, li2023large}, playing a crucial role in the development of data-augmented LLM applications. This approach not only extends the functional capabilities of LLMs but also enhances their practical utility across diverse domains.







\begin{figure}
    \centering
    \includegraphics[width=\textwidth]{contents/images/solution_conclusion4.png}
    \caption{Summary of Main Techniques for Different Query Levels in \MFUsentencecase{\app}}
    \label{fig:conclusion}
\end{figure}

\subsection{Fine-tuning}

Despite the robust in-context learning capabilities of LLMs, accurately identifying rationales or optimal examples for complex and lengthy logical chains remains a significant challenge. Additionally, the provision of extensive external prior knowledge can also pose challenges to the inference capabilities of LLMs. Given these factors, fine-tuning emerges as a promising approach. It not only utilizes the extensive foundational knowledge that LLMs acquire during pretraining but also enables them to rapidly grasp new domain rationales. This method provides a viable path for enhancing the adaptability and effectiveness of LLMs in tackling advanced and specialized tasks.


Instruction tuning is a common method for infusing new capabilities into LLMs, typically involving supervised fine-tuning using paired (instruction, output) data. There are three primary methods for constructing an instruction dataset: a) deriving from existing datasets~\cite{longpre2023flan, sanh2021multitask}, b) manually creating through handcrafted instructions~\cite{ouyang2022training, conover2023free, kopf2024openassistant}, and c) generating synthetic data using powerful LLMs~\cite{xu2023wizardlm, wang2022self}. Additionally, numerous studies~\cite{raffel2020exploring, iyer2022opt, wang2023far} have explored how to optimize the data distribution within instruction datasets to enhance fine-tuning effectiveness. However, when building data-augmented LLM applications, fine-tuning remains a relatively costly method in terms of time and computational resources. Recently, several efforts have been made to reduce the costs associated with fine-tuning large models. Adapter tuning, for instance, involves integrating small adapter models with LLMs while freezing the parameters of the LLM during fine-tuning and only optimizing the weights of the adapter~\cite{houlsby2019parameter, he2021towards, lei2023conditional, chronopoulou2023adaptersoup}. Prefix Tuning and Prompt Tuning involve adding a set of trainable vectors before the input, which are optimized during training to enhance the performance of the LLM~\cite{petrov2023prompting, mahabadi2021parameter, zhu2023spt, wang2023aprompt, choi2023smop}. Low-Rank Adaptation~\cite{hu2021lora, valipour2022dylora, zhang2023adalora, ding2023sparse, liu2024dora} reduces the number of trainable parameters needed for adapting to downstream tasks by imposing low-rank constraints on each dense layer to approximate the update matrices.


In recent years, there has been a substantial amount of work using supervised fine-tuning to enhance the capabilities of LLMs in specialized domains such as mathematical reasoning, finance, law, and healthcare~\cite{yu2022alert, zong2023solving, yang2024leandojo}. For instance, ChatTimeLlama~\cite{yuan2024back} introduced an interpretable time reasoning instruction tuning dataset and fine-tuned on LLaMA~\cite{touvron2023llama1} to significantly improve the model's complex temporal reasoning, future event prediction capabilities, and interpretability. LISA~\cite{lai2024lisa} leveraged a small set of segment data samples that involve reasoning to fine-tune the multimodal LLM LLaVA, which resulted in substantial improvements in reasoning segmentation capabilities. MAmmoTH~\cite{yue2023mammoth} ingeniously constructed a mathematical example dataset that uniquely combines Chain of Thought and Program of Thought reasoning, ensuring broad coverage across different mathematical domains and enhancing the LLM's ability to solve general mathematical problems. ReFT~\cite{trung2024reft} proposes a method for learning from multiple annotated reasoning paths corresponding to the same problem. It automatically samples numerous reasoning trajectories for a given mathematical problem, leveraging the correct answer to generate reward signals. ChatDoctor~\cite{li2023chatdoctor} utilized a large dataset of 100,000 patient-doctor dialogues from a widely-used online medical consultation platform to fine-tune LLaMA, significantly enhancing the model's ability to understand patient needs and provide effective recommendations. FinGPT~\cite{yang2023fingpt} developed an open-source LLM fine-tuned on financial data using automated data curation and lightweight, low-rank adaptation techniques. DISC-LawLLM~\cite{yue2023disc} created a supervised fine-tuning dataset for the Chinese judicial domain, fine-tuning LLMs to effectively serve a variety of users in different legal scenarios with enhanced legal reasoning capabilities.